% outlook/bim-DE.tex
\chapter{BIM und Alignment-Modellierung}

\section*{Motivation}

Building Information Modelling (BIM) ist zu einem zentralen Paradigma
moderner Infrastrukturprojekte geworden. Sein Ziel ist die Integration
von Geometrie, Semantik und Lebenszyklusinformationen in eine einheitliche
digitale Repräsentation.

Trotz dieses weiten Umfangs bleibt die Repräsentation von
Eisenbahn-Alignments in BIM begrifflich und technisch begrenzt.

\section{BIM als Paradigma der Datenintegration}

BIM befasst sich grundlegend mit:
\begin{itemize}
\item objektbasierter Modellierung,
\item semantischer Anreicherung von Geometrie,
\item fachübergreifender Interoperabilität,
\item Lebenszyklusmanagement von Infrastrukturanlagen.
\end{itemize}

Geometrie ist in diesem Kontext typischerweise Objekten wie Gebäuden,
Brücken, Gleiselementen und Geländemodellen zugeordnet.

Alignment-Geometrie ist häufig implizit in diese Objekte eingebettet,
anstatt als primäre Modellierungseinheit behandelt zu werden.

\section{Alignment in BIM-Standards}

\subsection{IFC Alignment}

Neuere Erweiterungen von IFC führen ausdrückliche Alignment-Entitäten ein.
Sie sollen Horizontalalignment, Vertikalalignment, Überhöhung und
Stationierung repräsentieren.

Obwohl dies einen bedeutenden Fortschritt darstellt, bleiben mehrere
Probleme bestehen:
\begin{itemize}
\item die Trennung der Komponenten,
\item das Fehlen einer einheitlichen Zustandsrepräsentation,
\item begrenzte Unterstützung dynamischer Interpretation,
\item die Komplexität der Schemanutzung.
\end{itemize}

\subsection{Geschichtete Repräsentation}

In den meisten BIM-bezogenen Formaten wird das Alignment in
Horizontalgeometrie, Vertikalprofil und Überhöhungsdefinition zerlegt.
Diese Zerlegung ist für den Datenaustausch praktisch, unterstützt jedoch
nicht unmittelbar Berechnung oder Optimierung.

\section{Grenzen gegenwärtiger BIM-Alignment-Modellierung}

\subsection{Statische Repräsentation}

BIM repräsentiert Geometrie hauptsächlich als statische Daten. Dynamische
Eigenschaften wie Beschleunigung, Ruck oder Fahrzeugantwort gehören nicht
zum Kernmodell.

\subsection{Fragmentierung}

Die Trennung horizontaler, vertikaler und überhöhungsbezogener Komponenten
führt zu einer fragmentierten Repräsentation. Diese Fragmentierung
erschwert Konsistenzprüfungen, dynamische Auswertung und
optimierungsbasierten Entwurf.

\subsection{Fehlender Berechnungskern}

BIM-Standards richten sich auf Datenaustausch und Dokumentation. Sie
stellen kein einheitliches Berechnungsmodell für Alignment-Analyse und
-Optimierung bereit.

\section{AIM als Berechnungsebene für BIM}

\subsection{Begriffliche Rolle}

Das AIM-Framework kann als Berechnungskern verstanden werden, der
BIM-Datenmodelle ergänzt. Statt BIM zu ersetzen, vermittelt AIM zwischen
externen Datenformaten und internem ingenieurwissenschaftlichem Schließen.

\subsection{Kanonische Repräsentation}

AIM führt eine einheitliche Repräsentation ein:
\[
(\gamma(s), \theta(s), \kappa(s), \phi(s)).
\]

Diese Repräsentation integriert Horizontalgeometrie sowie Vertikal- und
Überhöhungsverhalten, unterstützt dynamische Auswertung und ist direkt in
der Optimierung verwendbar.

\subsection{Transformationskette}

\begin{verbatim}
BIM / IFC / LandXML
→ Parsen und Quellevidenz
→ prüfbare Importkandidaten
→ ausdrückliche Zulassung oder Ablehnung
→ zugelassener Ingenieurgegenstand und qualifizierter Zustand
→ Analyse / Optimierung
→ Export
\end{verbatim}

Diese Kette trennt Belange des Datenaustauschs, die Verantwortung für
Evidenzprüfung und Zulassung, die Berechnungsmodellierung und das
ingenieurwissenschaftliche Schließen.

Erfolgreiches Parsen ist daher keine Zulassungsentscheidung. Importierte
Quelldatensätze können partiell, mehrdeutig oder ungeklärt bleiben; diese
Grenzen müssen mit ihrer Provenienz erhalten bleiben. Erst eine ausdrückliche
Zulassung kann einen dauerhaften Ingenieurgegenstand für nachgelagerte Services
verfügbar machen. Die gegenwärtige Referenzimplementierung ist eine
Demonstration dieser Trennung; sie bildet einen solchen Gegenstand als
SpotObject in ihrem geladenen SPOT-Universum ab, doch keiner der beiden
Implementierungsbegriffe definiert Kernel-Bedeutung oder Autorität. Dauerhafte
Unterscheidbarkeit bleibt dem aktiven Kernel-Datensatz KC-ID-001 Engineering
Object Identity zugeordnet.

\section{Alignment-zentriertes BIM}

\subsection{Perspektivwechsel}

Das AIM-Framework legt einen Wechsel zur Alignment-zentrierten
Modellierung nahe. Statt das Alignment als eine Komponente unter vielen
zu behandeln, wird es zur primären Ordnungsstruktur.

\subsection{Folgerungen}

\begin{itemize}
\item Infrastrukturelemente werden auf das Alignment bezogen,
\item Geometrie wird aus dem Alignment-Zustand abgeleitet,
\item dynamisches Verhalten wird Teil des Modells,
\item Optimierung wird zu einer nativen Operation.
\end{itemize}

\section{Integration in BIM-Arbeitsabläufe}

\subsection{Import}

Bestehende BIM-Daten werden durch IFC- und LandXML-Parser sowie durch die
Interpretation von Altformaten untersucht.

\subsection{Verarbeitung}

Innerhalb von AIM dürfen zugelassene Informationen normalisiert,
rekonstruiert, um eine dynamische Interpretation angereichert und
gegebenenfalls optimiert werden.

Nicht zugelassene Informationen bleiben Import-Evidenz. Fehlende
Objektrelationen, Koordinatenbedeutung, Kilometrierungsinterpretation oder
Quellenautorität dürfen nicht aus Dateinamen, Importreihenfolge oder
repräsentationaler Ähnlichkeit abgeleitet werden. Mehrere synchronisierte
Workspace-Sichten können ein zugelassenes Alignment darstellen, bleiben aber
Repräsentationen desselben Objekts und derselben Revision.

\subsection{Export}

Ergebnisse können in IFC-Alignment-Strukturen, LandXML-Dateien oder andere
Ingenieurformate zurückexportiert werden. Dies gewährleistet die
Kompatibilität mit bestehenden BIM-Ökosystemen.

\section{Bezug zu BIMinfra}

Infrastruktur-BIM (häufig als BIMinfra bezeichnet) erweitert BIM-Begriffe
auf lineare Infrastruktur wie Eisenbahnen und Straßen.

AIM kann als Spezialisierung innerhalb von BIMinfra verstanden werden,
die sich auf Alignment-zentrierte Modellierung, dynamische Interpretation
und optimierungsgetriebenen Entwurf richtet.

\section{Diskussion}

BIM stellt ein leistungsfähiges Framework zur Datenintegration bereit,
besitzt jedoch kein starkes Berechnungsmodell für Alignments.

AIM schließt diese Lücke durch eine einheitliche Zustandsrepräsentation,
dynamische Auswertung, Optimierungsunterstützung und die Integration
heterogener Datenquellen.

BIM und AIM können somit als komplementär gelten: BIM für Datenintegration
und Lebenszyklus, AIM für Berechnung und Alignment-bezogenes Schließen.

\section{Zusammenfassung}
\begin{summary}
Gegenwärtige BIM-Standards stellen wichtige Strukturen für Datenaustausch
und semantische Modellierung bereit, bleiben jedoch in ihrer Behandlung
des Eisenbahn-Alignments begrenzt.

Das AIM-Framework ergänzt BIM durch eine einheitliche, zustandsbasierte
und optimierungsbereite Repräsentation des Alignments.

Dies ermöglicht den Übergang von statischer Datenverarbeitung zu
dynamischer, berechnungsorientierter und Alignment-zentrierter
Infrastrukturmodellierung.
\end{summary}
